\subsection{Parameter Recovery under Nonlinear Effort}\label{app-effort-recovery}

{\color{blue}To assess the effort approximation, we generate data under nonlinear equilibrium policies and estimate parameters with the leading-order likelihood. Ten joint parameter configurations represent high-incentive contests in the original estimates, with four realizations per configuration. All six parameters are re-estimated using the empirical priors in Section~\ref{sec-kaggle-application}. The two exploratory realizations used in prior development are excluded. Table~\ref{tab-nonlinear-recovery} reports recovery of feedback precision and innovation risk.}

\begin{table}[!htbp]
\centering
\caption{Parameter Recovery from Nonlinear Simulated Contests}
\label{tab-nonlinear-recovery}
{\small
\begin{tabularx}{\textwidth}{@{}l>{\centering\arraybackslash}X>{\centering\arraybackslash}X@{}}
\toprule
Parameter & \makecell{Mean absolute\\relative error (\%)} & \makecell{Coverage of 95\%\\credible intervals (\%)}\\
\midrule
Signal precision $\lambda$ & 17.0 & 97.5\\
Innovation risk $\sigma$ & 12.1 & 82.5\\
\bottomrule
\end{tabularx}
\par\vspace{1ex}\noindent
\begin{minipage}{\textwidth}
\footnotesize\textit{Note:} Each row summarizes 40 simulated datasets at ten fixed parameter configurations. Point estimates are posterior medians. Absolute relative error divides the absolute estimation error by the generating value. Coverage is the fraction of datasets whose equal-tailed 95\% posterior interval contains that value. The data are generated under nonlinear-equilibrium effort and estimated with the leading-order likelihood; each of the ten fixed configurations, chosen to represent high-incentive empirical contests, contributes four realizations. {\color{blue}Thus, the entries are equally weighted summaries of these 40 datasets, not estimates from a random sample of all contests. Lower relative error means more accurate point recovery; coverage assesses interval calibration. Innovation-risk coverage of 82.5\% is below the nominal 95\%, and precision coverage of 97.5\% does not establish generally calibrated inference.}
\end{minipage}}
\end{table}

For feedback precision, the 95\% credible interval contains the generating value in 14 of the 15 realizations with point-estimation errors above 20\%; only one of the 40 datasets has an interval that misses the generating value. For innovation risk, coverage is 33 of 40 datasets (82.5\%), below the nominal 95\% level. This flags possible undercoverage in the tested configurations; the small, heterogeneous simulation sample does not establish its magnitude more generally or isolate its cause.

{\color{blue}The simulation and its likelihood use interval-start public feedback in the perceived-state recursion; the empirical application uses the interval-end public value after computing that interval's effort. Both likelihoods score fresh observations at interval endpoints. The recovery results therefore concern these ten configurations and the simulator-aligned feedback convention, rather than providing an error bound for all empirical estimates.}

